% ---------------------------------------------------------------- 3.6 B3: DPO in practice
\begin{frame}{DPO in Practice: What You Give Up}
DPO's simplicity has a cost: you lose the control knobs the RL loop gave you.
\begin{itemize}\itemsep0.5em\small
    \item \textbf{It constrains only the \emph{gap}, nothing else.} The loss just
          pushes the preferred answer's favoring above the rejected one's. With no
          on-policy sampling to pin down absolute behaviour, it can push \emph{both}
          down, leaking probability onto unrelated answers it never saw. RLHF's
          RL loop and ``don't-drift'' anchor would catch that; DPO has neither.
    \item \textbf{Length bias (shared, but harder to patch).} If the data favours
          long answers, \emph{both} RLHF and DPO turn verbose. RLHF lets you subtract
          an explicit length penalty from the reward; DPO has no separate reward to
          adjust, so length control must be bolted on as extra loss terms.
    \item \textbf{No fresh samples.} The policy only ever sees the fixed preference
          pairs, so it can drift off that distribution, the price of dropping the
          RL loop.
\end{itemize}
\end{frame}

